\chapter{Standard Functions}\label{ch:b1:functions}

Analysis is only as useful as the stock of functions one masters. To
the collection inherited from secondary school --- powers, exponential,
logarithm, trigonometric functions --- this chapter adds their inverse
functions ($\arcsin$, $\arccos$, $\arctan$) and the hyperbolic family.
Derivatives are used freely at the level of \cref{ch:g12:deriv}; the
theory behind inverse functions is completed in
\cref{ch:b1:continuity,ch:b1:derivative}.

\section{Exponential, logarithm, powers}

\begin{proposition}[Recap and characterization]\label{prop:b1:functions:expln}
$\exp \colon \R \to \intoo{0}{+\infty}$ and $\ln \colon
\intoo{0}{+\infty} \to \R$ are reciprocal bijections, strictly
increasing, with
\[
\exp(x + y) = \exp x \exp y, \qquad \ln(xy) = \ln x + \ln y,
\qquad \exp' = \exp, \qquad \ln'(x) = \frac 1x .
\]
Moreover $\exp$ is the \emph{only} differentiable function $f \colon
\R \to \R$ with $f' = f$ and $f(0) = 1$.
\end{proposition}

\begin{proof}
The recap is \cref{ch:g12:exp}. For uniqueness, let $f' = f$,
$f(0) = 1$, and set $g(x) = f(x)\,\eu^{-x}$. Then $g' = f'\eu^{-x} -
f\eu^{-x} = 0$, so $g$ is constant equal to $g(0) = 1$: $f = \exp$.
\end{proof}

\begin{definition}[General powers]\label{def:b1:functions:powers}
For $x > 0$ and $\alpha \in \R$:
$\;x^\alpha = \eu^{\alpha \ln x}$\index{power function}. For $a > 0$,
$a \neq 1$, the \emph{logarithm to base $a$} is $\log_a x =
\frac{\ln x}{\ln a}$, the inverse of $x \mapsto a^x$.
\end{definition}

\begin{proposition}[Power rules and growth comparison]\label{prop:b1:functions:powerrules}
For $x, y > 0$ and $\alpha, \beta \in \R$:
\[
x^{\alpha+\beta} = x^\alpha x^\beta, \quad
(x^\alpha)^\beta = x^{\alpha\beta}, \quad
(xy)^\alpha = x^\alpha y^\alpha, \quad
(x^\alpha)' = \alpha\, x^{\alpha - 1}.
\]
Growth scale\index{growth comparison} as $x \to +\infty$, for every
$\alpha > 0$ and $\beta > 0$:
\[
\frac{(\ln x)^\beta}{x^\alpha} \longrightarrow 0,
\qquad
\frac{x^\alpha}{\eu^{\beta x}} \longrightarrow 0 .
\]
\end{proposition}

\begin{proof}
The identities transcribe those of $\exp$ and $\ln$ through the
definition; the derivative is the chain rule: $(\eu^{\alpha\ln x})' =
\frac{\alpha}{x} \eu^{\alpha\ln x}$.

Comparisons: from $\frac{\ln t}{t} \to 0$ (proved in
\cref{ch:g12:exp}), substitute $t = x^{\alpha/\beta}$:
$\frac{\ln x}{x^{\alpha/\beta}} \to 0$, then raise to the power
$\beta$. For the second, substitute $x = \ln u$ in the first.
\end{proof}

\section{Inverse trigonometric functions}

\begin{definition}[$\arcsin$, $\arccos$, $\arctan$]\label{def:b1:functions:arc}
The restrictions
\[
\sin \colon \intcc{-\tfrac\pi2}{\tfrac\pi2} \to \intcc{-1}{1},
\qquad
\cos \colon \intcc{0}{\pi} \to \intcc{-1}{1},
\qquad
\tan \colon \intoo{-\tfrac\pi2}{\tfrac\pi2} \to \R
\]
are strictly monotonic bijections. Their inverse maps are written
$\arcsin$\index{arcsine}, $\arccos$\index{arccosine} and
$\arctan$\index{arctangent}. Thus, for instance, $y = \arcsin x$
($x \in \intcc{-1}{1}$) is \emph{the} angle in
$\intcc{-\frac\pi2}{\frac\pi2}$ whose sine is $x$.
\end{definition}

\begin{omfigure}
\begin{tikzpicture}
\begin{axis}[omaxis, width=5.6cm, height=5.6cm,
    xmin=-1.25, xmax=1.25, ymin=-1.8, ymax=1.8,
    xtick={-1,1}, ytick={-1.5708,1.5708},
    yticklabels={$-\frac{\pi}{2}$,$\frac{\pi}{2}$}]
  \addplot[omDef, very thick, domain=-1:1, samples=200]
    {rad(asin(x))};
\end{axis}
\begin{scope}[xshift=5.3cm]
\begin{axis}[omaxis, width=5.6cm, height=5.6cm,
    xmin=-1.25, xmax=1.25, ymin=-0.4, ymax=3.4,
    xtick={-1,1}, ytick={1.5708,3.1416},
    yticklabels={$\frac{\pi}{2}$,$\pi$}]
  \addplot[omThm, very thick, domain=-1:1, samples=200]
    {rad(acos(x))};
\end{axis}
\end{scope}
\begin{scope}[xshift=10.6cm]
\begin{axis}[omaxis, width=5.6cm, height=5.6cm,
    xmin=-4, xmax=4, ymin=-1.9, ymax=1.9,
    xtick={-2,2}, ytick={-1.5708,1.5708},
    yticklabels={$-\frac{\pi}{2}$,$\frac{\pi}{2}$}]
  \addplot[omMeth, very thick, domain=-4:4, samples=200]
    {rad(atan(x))};
  \addplot[dashed, gray, domain=-4:4] {1.5708};
  \addplot[dashed, gray, domain=-4:4] {-1.5708};
\end{axis}
\end{scope}
\end{tikzpicture}

{\small From left to right: $\arcsin$, $\arccos$, $\arctan$. Each
inherits its graph from the restricted direct function by reflection
in the line $y = x$.}
\end{omfigure}

\begin{proposition}[Derivatives]\label{prop:b1:functions:arcderiv}
On the interior of their domains:
\[
\arcsin' x = \frac{1}{\sqrt{1 - x^2}},
\qquad
\arccos' x = \frac{-1}{\sqrt{1 - x^2}},
\qquad
\arctan' x = \frac{1}{1 + x^2}.
\]
\end{proposition}

\begin{proof}
Anticipating the inverse-function rule proved in
\cref{ch:b1:derivative}: if $f$ is a differentiable bijection and
$f' \neq 0$, then $(f^{-1})'(x) = \frac{1}{f'(f^{-1}(x))}$. For
$\arcsin$: with $y = \arcsin x$,
\[
\arcsin' x = \frac{1}{\cos y}
= \frac{1}{\sqrt{1 - \sin^2 y}} = \frac{1}{\sqrt{1 - x^2}},
\]
where $\cos y = +\sqrt{1 - \sin^2 y}$ because $y \in
\intoo{-\frac\pi2}{\frac\pi2}$ forces $\cos y > 0$. Similarly
$\arccos' x = \frac{1}{-\sin y}$ with $\sin y > 0$ on
$\intoo{0}{\pi}$, and $\arctan' x = \frac{1}{1 + \tan^2 y} =
\frac{1}{1 + x^2}$ using $\tan' = 1 + \tan^2$.
\end{proof}

\begin{proposition}[Standard identities]\label{prop:b1:functions:arcidentities}
\begin{enumerate}
  \item For $x \in \intcc{-1}{1}$: $\arcsin x + \arccos x =
        \dfrac{\pi}{2}$.
  \item For $x > 0$: $\arctan x + \arctan\dfrac 1x = \dfrac{\pi}{2}$
        (and $-\frac\pi2$ for $x < 0$).
  \item $\sin(\arccos x) = \cos(\arcsin x) = \sqrt{1 - x^2}$;
        $\;\tan(\arcsin x) = \dfrac{x}{\sqrt{1 - x^2}}$ for
        $\abs x < 1$.
\end{enumerate}
\end{proposition}

\begin{proof}
(1) The derivative of $x \mapsto \arcsin x + \arccos x$ is zero on
$\intoo{-1}{1}$ (\cref{prop:b1:functions:arcderiv}), so the function
is constant there, equal to its value $\frac\pi2$ at $0$; the endpoint
values $\pm 1$ are checked directly ($\frac\pi2 + 0$ and $-\frac\pi2 +
\pi$).

(2) Same method on $\intoo{0}{+\infty}$: the derivative is
$\frac{1}{1+x^2} + \frac{-1/x^2}{1 + 1/x^2} = \frac{1}{1+x^2} -
\frac{1}{x^2 + 1} = 0$, and at $x = 1$ the sum is $2\arctan 1 =
\frac\pi2$. For $x < 0$, use oddness of $\arctan$.

(3) With $y = \arccos x \in \intcc{0}{\pi}$: $\sin y \geq 0$, so
$\sin y = \sqrt{1 - \cos^2 y} = \sqrt{1 - x^2}$; likewise for
$\cos(\arcsin x)$, and the tangent formula is the quotient.
\end{proof}

\begin{remark}
$\arcsin(\sin\theta) = \theta$ holds \emph{only} for $\theta \in
\intcc{-\frac\pi2}{\frac\pi2}$: for instance
$\arcsin(\sin\pi) = 0 \neq \pi$. The composition in the other
direction, $\sin(\arcsin x) = x$, is valid on all of
$\intcc{-1}{1}$.
\end{remark}

\section{Hyperbolic functions}

\begin{definition}[$\cosh$, $\sinh$, $\tanh$]\label{def:b1:functions:hyperbolic}
For $x \in \R$:
\[
\cosh x = \frac{\eu^x + \eu^{-x}}{2},
\qquad
\sinh x = \frac{\eu^x - \eu^{-x}}{2},
\qquad
\tanh x = \frac{\sinh x}{\cosh x}
\]
(\emph{hyperbolic cosine, sine, tangent}\index{hyperbolic
functions}). $\cosh$ is even, $\sinh$ and $\tanh$ are odd.
\end{definition}

\begin{proposition}[Basic properties]\label{prop:b1:functions:hyprules}
For all $x, y \in \R$:
\begin{enumerate}
  \item $\cosh^2 x - \sinh^2 x = 1$;
  \item $\cosh' = \sinh$, $\;\sinh' = \cosh$, $\;\tanh' = 1 - \tanh^2
        = \dfrac{1}{\cosh^2}$;
  \item $\cosh(x+y) = \cosh x\cosh y + \sinh x \sinh y$ and
        $\sinh(x+y) = \sinh x \cosh y + \cosh x \sinh y$;
  \item $\sinh$ is a strictly increasing bijection of $\R$ onto $\R$;
        $\cosh$ restricted to $\R_+$ is a strictly increasing
        bijection onto $\intco{1}{+\infty}$; $\tanh$ is a strictly
        increasing bijection of $\R$ onto $\intoo{-1}{1}$.
\end{enumerate}
\end{proposition}

\begin{proof}
(1) $\bigl(\frac{\eu^x + \eu^{-x}}{2}\bigr)^2 - \bigl(\frac{\eu^x -
\eu^{-x}}{2}\bigr)^2 = \frac{(\eu^{2x} + 2 + \eu^{-2x}) - (\eu^{2x} -
2 + \eu^{-2x})}{4} = 1$.

(2) Differentiate the defining formulas; for $\tanh$, the quotient
rule gives $\frac{\cosh^2 - \sinh^2}{\cosh^2}$, which is both
$1 - \tanh^2$ and $\frac{1}{\cosh^2}$ by (1).

(3) Expand the right-hand sides using the definitions; the cross terms
assemble into $\frac{\eu^{x+y} + \eu^{-(x+y)}}{2}$, resp.\
$\frac{\eu^{x+y} - \eu^{-(x+y)}}{2}$.

(4) $\sinh' = \cosh \geq 1 > 0$, so $\sinh$ is strictly increasing,
with limits $\pm\infty$ (dominant term $\pm\frac{\eu^{\abs
x}}{2}$); the bijection statement then follows from the intermediate
value theorem (\cref{ch:g12:limcont}; systematic treatment in
\cref{ch:b1:continuity}). On $\R_+$, $\cosh' = \sinh \geq 0$
vanishing only at $0$: strictly increasing from $\cosh 0 = 1$ to
$+\infty$. And $\tanh' > 0$ with limits $\pm 1$ at $\pm\infty$
(divide numerator and denominator by $\eu^x$).
\end{proof}

\begin{omfigure}
\begin{tikzpicture}
\begin{axis}[omaxis, width=6.4cm, height=6cm,
    xmin=-2.6, xmax=2.6, ymin=-3.6, ymax=3.9,
    xtick={-2,-1,1,2}, ytick={-3,-2,-1,1,2,3}]
  \addplot[omDef, very thick, domain=-2.05:2.05, samples=120]
    {cosh(x)} node[pos=0.1, right, font=\small] {$\cosh$};
  \addplot[omThm, very thick, domain=-2.05:2.05, samples=120]
    {sinh(x)} node[pos=0.93, above left, font=\small] {$\sinh$};
  \addplot[dashed, gray, domain=0:2.5] {exp(x)/2};
\end{axis}
\begin{scope}[xshift=7.4cm]
\begin{axis}[omaxis, width=6.4cm, height=6cm,
    xmin=-3.2, xmax=3.2, ymin=-1.5, ymax=1.5,
    xtick={-2,-1,1,2}, ytick={-1,1}]
  \addplot[omMeth, very thick, domain=-3:3, samples=150]
    {tanh(x)} node[pos=0.9, above left, font=\small] {$\tanh$};
  \addplot[dashed, gray, domain=-3:3] {1};
  \addplot[dashed, gray, domain=-3:3] {-1};
\end{axis}
\end{scope}
\end{tikzpicture}

{\small Left: $\cosh$ and $\sinh$, asymptotically glued to
$\frac{\eu^x}{2}$ (dashed). Right: $\tanh$, increasing from $-1$ to
$1$.}
\end{omfigure}

\begin{remark}[Why ``hyperbolic''?]
The point $(\cosh t, \sinh t)$ runs along the branch $x > 0$ of the
hyperbola $x^2 - y^2 = 1$ (by
\cref{prop:b1:functions:hyprules}~(1)), exactly as $(\cos t, \sin
t)$ runs along the circle $x^2 + y^2 = 1$. Every circular identity has
a hyperbolic sibling, with sign changes governed by
$\sinh^2 \leftrightarrow -\sin^2$.
\end{remark}

\begin{proposition}[Inverse hyperbolic functions]\label{prop:b1:functions:invhyp}
The inverses granted by \cref{prop:b1:functions:hyprules}~(4) have
closed forms:
\[
\operatorname{arsinh} x = \ln\bigl(x + \sqrt{x^2 + 1}\bigr)
\ (x \in \R),
\qquad
\operatorname{arcosh} x = \ln\bigl(x + \sqrt{x^2 - 1}\bigr)
\ (x \geq 1),
\]
\[
\operatorname{artanh} x = \frac 12 \ln\frac{1 + x}{1 - x}
\ (\abs x < 1),
\]
with derivatives $\frac{1}{\sqrt{x^2+1}}$, $\frac{1}{\sqrt{x^2-1}}$
($x > 1$) and $\frac{1}{1 - x^2}$ respectively.
\end{proposition}

\begin{proof}
For $\operatorname{arsinh}$: solve $x = \sinh y = \frac{\eu^y -
\eu^{-y}}{2}$. Setting $u = \eu^y > 0$: $u^2 - 2xu - 1 = 0$, so $u = x
+ \sqrt{x^2 + 1}$ (the root $x - \sqrt{x^2+1}$ is negative), and $y =
\ln(x + \sqrt{x^2+1})$. The other two are identical computations
($u^2 - 2xu + 1 = 0$ for $\operatorname{arcosh}$, keeping the root
$\geq 1$; a two-line rearrangement for $\operatorname{artanh}$).
Derivatives: differentiate the logarithmic expressions, e.g.
\[
\bigl(\ln(x + \sqrt{x^2+1})\bigr)'
= \frac{1 + \frac{x}{\sqrt{x^2+1}}}{x + \sqrt{x^2+1}}
= \frac{1}{\sqrt{x^2 + 1}} . \qedhere
\]
\end{proof}

\begin{method}[Choosing the right primitive form]\label{met:b1:functions:primitive}
The three derivative patterns to memorize for integration
(\cref{ch:b1:integration}):
\[
\int \frac{\dd x}{1 + x^2} = \arctan x + C,
\qquad
\int \frac{\dd x}{\sqrt{1 - x^2}} = \arcsin x + C,
\qquad
\int \frac{\dd x}{\sqrt{x^2 + 1}} = \operatorname{arsinh} x + C,
\]
and for a general quadratic, reduce to these by completing the square
and rescaling.
\end{method}

\section{Exercises}

\begin{exercise}[$\star$]\label{exo:b1:functions:1}
Compute without calculator: $\arcsin\frac{\sqrt 3}{2}$;
$\;\arccos\bigl(-\frac 12\bigr)$; $\;\arctan(-1)$;
$\;\arcsin\bigl(\sin\frac{5\pi}{6}\bigr)$;
$\;\arccos\bigl(\cos\frac{7\pi}{4}\bigr)$.
\end{exercise}

\begin{exercise}[$\star$]\label{exo:b1:functions:2}
Give the domain of definition of $f(x) = \arcsin(2x - 1)$ and compute
$f'$ where defined. Same questions for $g(x) = \arctan\sqrt{x}$.
\end{exercise}

\begin{exercise}[$\star$]\label{exo:b1:functions:3}
Prove that for all $x \in \R$: $\cosh x + \sinh x = \eu^x$ and
$(\cosh x + \sinh x)^n = \cosh nx + \sinh nx$ for every $n \in \Z$
(a hyperbolic de Moivre formula).
\end{exercise}

\begin{exercise}[$\star$]\label{exo:b1:functions:4}
Simplify $\cos(2\arcsin x)$ and $\sin(2\arctan x)$ into algebraic
expressions of $x$.
\end{exercise}

\begin{exercise}[$\star$]\label{exo:b1:functions:5}
Order, for large $x$, the functions
$x^{100}$, $\;\eu^{\sqrt x}$, $\;(\ln x)^{1000}$, $\;\eu^{x/100}$,
$\;x^{\ln x}$, from slowest to fastest growth, with justifications
based on \cref{prop:b1:functions:powerrules}.
\end{exercise}

\begin{exercise}[$\star\star$]\label{exo:b1:functions:6}
Study the function $f(x) = \arctan\dfrac{2x}{1 - x^2}$ on its domain:
compute $f'$, compare with $(2\arctan x)'$, and express $f(x)$ in
terms of $\arctan x$ on each of the three intervals of the domain.
\end{exercise}

\begin{exercise}[$\star\star$]\label{exo:b1:functions:7}
Solve in $\R$: $\;\cosh x = 2$; then $5\cosh x - 4 \sinh x = 3$
(express the solutions with logarithms). \emph{Hint for the second:
write everything with $u = \eu^x$.}
\end{exercise}

\begin{exercise}[$\star\star$]\label{exo:b1:functions:8}
Prove the identity $\arctan\frac{1}{2} + \arctan\frac{1}{3} =
\frac{\pi}{4}$, then Machin's formula\index{Machin's formula}
\[
4\arctan\frac 15 - \arctan\frac{1}{239} = \frac{\pi}{4}.
\]
\emph{Hint: compute the tangent of both sides using the addition
formula, and control in which interval each side lies.}
\end{exercise}

\begin{exercise}[$\star\star$]\label{exo:b1:functions:9}
Prove that for all $x \geq 0$: $\;x - \dfrac{x^3}{6} \leq \sin x \leq
x$ \emph{(study the successive derivatives of the differences)}, and
deduce $\lim_{x \to 0^+} \frac{x - \sin x}{x^3} = \frac 16$ is
plausible --- the limit itself is established in
\cref{ch:b1:taylor}.
\end{exercise}

\begin{exercise}[$\star\star\star$]\label{exo:b1:functions:10}
For $x \in \intcc{-1}{1}$, set
$h(x) = \arcsin\bigl(2x\sqrt{1 - x^2}\,\bigr) - 2\arcsin x$. One might
expect $h = 0$ from the identity $\sin 2y = 2\sin y\cos y$ --- but
$h$ is \emph{not} identically zero. Compute $h'$ on the open intervals
where it exists, evaluate $h$ at well-chosen points, and give the full
piecewise-constant description of $h$ on $\intcc{-1}{1}$.
\end{exercise}

\begin{exercise}[$\star\star\star$]\label{exo:b1:functions:11}
(Gudermannian) Let $g(x) = \arctan(\sinh x)$ for $x \in \R$. Prove
that $g$ is an odd, strictly increasing bijection from $\R$ onto
$\intoo{-\frac\pi2}{\frac\pi2}$, that $g'(x) = \frac{1}{\cosh x}$,
and that $\tan g(x) = \sinh x$, $\;\sin g(x) = \tanh x$,
$\;\cos g(x) = \frac{1}{\cosh x}$: the function $g$ links circular
and hyperbolic trigonometry without complex numbers.
\end{exercise}
